%% Relatório final — Projeto LLM 2026/2 — Grupo 1
%% Template: ACM sigconf (acmart), com os metadados de exemplo da ACM removidos.

% Arquivo .bib gerado automaticamente a partir deste .tex (dispensa um .bib separado).
\begin{filecontents*}[overwrite]{referencias.bib}
@article{schmidgall2025agentlab,
  title   = {Agent Laboratory: Using {LLM} Agents as Research Assistants},
  author  = {Schmidgall, Samuel and Su, Yusheng and Wang, Ze and Sun, Ximeng
             and Wu, Jialian and Yu, Xiaodong and Liu, Jiang and Moor, Michael
             and Liu, Zicheng and Barsoum, Emad},
  journal = {arXiv preprint arXiv:2501.04227},
  year    = {2025}
}
\end{filecontents*}

\documentclass[sigconf]{acmart}
\usepackage[brazil]{babel}

\AtBeginDocument{%
  \providecommand\BibTeX{{%
    Bib\TeX}}}

%% Remove o bloco de copyright/conferência e a "ACM Reference Format",
%% que não se aplicam a um relatório de disciplina.
\setcopyright{none}
\settopmatter{printacmref=false}
\renewcommand\footnotetextcopyrightpermission[1]{}
\pagestyle{plain}

\begin{document}

\title{Memória de Lições entre Projetos no Agent Laboratory}
\subtitle{Relatório final --- Projeto LLM 2026/2, Grupo 1}

\author{Déborah Yamamoto}
\affiliation{%
  \institution{DCC -- Universidade Federal de Minas Gerais}
  \city{Belo Horizonte}
  \state{MG}
  \country{Brasil}}

\author{Arthur SOBRENOME} % TODO: sobrenome do Arthur
\affiliation{%
  \institution{DCC -- Universidade Federal de Minas Gerais}
  \city{Belo Horizonte}
  \state{MG}
  \country{Brasil}}

\author{Cecília Junqueira}
\affiliation{%
  \institution{DCC -- Universidade Federal de Minas Gerais}
  \city{Belo Horizonte}
  \state{MG}
  \country{Brasil}}

\renewcommand{\shortauthors}{Yamamoto, SOBRENOME e Junqueira}

% Rascunho de resumo — completar com os resultados.
\begin{abstract}
O Agent Laboratory é um sistema de múltiplos agentes LLM que conduz
autonomamente um projeto de pesquisa, da revisão de literatura ao relatório,
mas descarta ao fim de cada projeto todo o feedback recebido. Propomos uma
memória de lições que sobrevive entre projetos: o feedback é condensado em
lições do tipo ``o que deu errado e como evitar'' e reaproveitado nos projetos
seguintes. Comparamos o sistema original com duas versões com memória, uma
alimentada pelo revisor LLM e outra por feedback humano, para avaliar se
guardar lições ajuda e se a origem da lição importa.
% TODO: uma ou duas frases com os principais resultados.
\end{abstract}

\keywords{agentes LLM, memória, feedback humano, pesquisa automatizada}

\maketitle

\section{Problema e Motivação}

O Agent Laboratory~\cite{schmidgall2025agentlab} é um sistema de múltiplos
agentes LLM que recebe uma ideia de pesquisa e conduz, de forma autônoma, a
revisão de literatura, a execução de experimentos e a redação do relatório.
Ao longo de cada projeto, o sistema recebe feedback de revisores LLM e,
opcionalmente, de um humano. Ao fim do projeto, porém, tudo o que foi
aprendido é descartado: cada nova pesquisa começa do zero.

Três evidências do próprio artigo motivam este trabalho. Primeiro, o
revisor automático é otimista demais: ele atribui, em média, nota 6,1 (de 10)
aos artigos gerados, contra 3,8 dada por revisores humanos. Segundo, o
feedback humano a cada etapa (modo \emph{co-pilot}) eleva a nota geral em
0,58 ponto, o que indica que o feedback certo melhora o agente. Terceiro, os
autores relatam falhas que se repetem entre projetos, como buscas em
\emph{loop} e experimentos que terminam com 0\% de acurácia. Se os erros se
repetem e o feedback ajuda, mas o aprendizado não é guardado, há uma
oportunidade clara de reaproveitar lições de um projeto no seguinte.

Investigamos duas questões de pesquisa:
\begin{itemize}
  \item \textbf{QP1.} Guardar lições de projetos anteriores faz o agente se
        sair melhor nos projetos seguintes?
  \item \textbf{QP2.} Importa a origem da lição: o revisor automático ou um
        humano?
\end{itemize}

\section{Abordagem}

Mantemos o fluxo e o código aberto do Agent Laboratory, composto de três
fases: (i)~\emph{literatura}, que busca e resume artigos;
(ii)~\emph{experimentos}, que planeja, codifica e executa; e
(iii)~\emph{relatório}, que escreve os resultados. Ao longo das fases, o
sistema recebe \emph{feedback} de um revisor LLM ou de um humano.
A única alteração é a adição de uma \textbf{memória de lições} que sobrevive
entre projetos.

Ao fim de cada projeto, o feedback recebido é condensado em lições no formato
``o que deu errado e como evitar'' e gravado na memória. Quando um novo
projeto começa, o agente consulta essa memória e passa a contar com as lições
acumuladas nas etapas seguintes. Assim, a memória fecha o ciclo que o sistema
original deixa aberto: o feedback, que antes só servia ao projeto corrente,
passa a influenciar os próximos.
% TODO: se couber, descrever em uma frase como as lições são recuperadas e
% em quais prompts/etapas são inseridas.

Para responder às questões de pesquisa, comparamos três versões do agente,
que diferem apenas na memória e na origem das lições:
\begin{itemize}
  \item \textbf{A (baseline):} Agent Laboratory original, sem memória;
  \item \textbf{B:} memória com lições extraídas do feedback do revisor LLM;
  \item \textbf{C:} memória com lições extraídas do feedback humano (do grupo).
\end{itemize}
A comparação entre A e as versões com memória responde à QP1; a comparação
entre B e C responde à QP2. Nossa hipótese é C~$>$~B~$>$~A. Consideramos
plausível, porém, que B piore o agente em relação a A: como o revisor
automático é otimista demais, as lições derivadas dele podem ser pouco
informativas ou até enganosas, o que também seria um resultado relevante.

\bibliographystyle{ACM-Reference-Format}
\bibliography{referencias}

\end{document}